\documentclass[a4paper]{article}
\usepackage{titlesec}
\usepackage{geometry}
\geometry{hmargin={2cm,2cm}, vmargin={3.5cm,3.5cm}}
\usepackage{graphicx}
\usepackage[table]{xcolor}
\usepackage{hyperref}
\hypersetup{pdfborder = {0 0 0},
colorlinks,
citecolor=black,
filecolor=black,
linkcolor=black,
urlcolor=blue}
\definecolor{logo}{RGB}{20,99,163}
\usepackage[T1]{fontenc}
\usepackage[utf8]{inputenc}
\usepackage[english,italian]{babel}
\usepackage{fancyhdr}
\usepackage{lastpage}
\usepackage{comment}

% creazione comando subsubsubsection
\titleclass{\subsubsubsection}{straight}[\subsection]

\newcounter{subsubsubsection}[subsubsection]
\renewcommand\thesubsubsubsection{\thesubsubsection.\arabic{subsubsubsection}}
\titleformat{\subsubsubsection}
    {\normalfont\normalsize\bfseries}{\thesubsubsubsection}{1em}{}
    \titlespacing*{\subsubsubsection}
    {0pt}{3.25ex plus 1ex minus .2ex}{1.5ex plus .2ex}

\makeatletter
\def\toclevel@subsubsubsection{4}
\def\l@subsubsubsection{\@dottedtocline{4}{7em}{4em}}
\makeatother

\setcounter{secnumdepth}{4}
\setcounter{tocdepth}{4}
%fine creazione comando


%Comandi di impaginazione uguale per tutti i documenti
\pagestyle{fancy}
\lhead{\includegraphics[scale=0.2]{../../../../Images/logo.png}}
%Titolo del documento
\rhead{\large Verbale Interno del 2021-05-10}
%\rfoot{\thepage}
\cfoot{Pagina \thepage{} di \pageref{LastPage}}
\setlength{\headheight}{45pt}
\renewcommand{\footrulewidth}{0.4pt}
\begin{document}

\begin{titlepage}
    \begin{center}
        \includegraphics{../../../../Images/logo.png}\\
        \vspace{20px}
        \textcolor{logo}{\hrulefill}\\
        \vspace{20px}
        \textbf{\huge\textcolor{logo}{Verbale Interno del 2021-05-10}}\\
        \vspace{10px}
        \textcolor{logo}{\hrulefill}\\
        \vspace{40px}
        \textbf{\Large Informazioni sul documento}\\
        \vspace{20px}
        \begin{tabular}{p{100px} | p{100px}}
            \textbf{Versione}     & 1.0.0                                                                      \\
            \textbf{Responsabile} & Tito Scutari                                                               \\
            \textbf{Redattore}    & Marco Barbaresco                                                           \\
            \textbf{Verificatore} & Nicolò Giaccone                                                            \\
            \textbf{Destinatari}  & Prof. Tullio Vardanega \newline Prof. Riccardo Cardin, \newline TechSWEave \\
            \textbf{Stato}        & Approvato                                                                  \\
            \textbf{Uso}          & Interno                                                                    \\
        \end{tabular}\\
        \vspace{60px}
        \href{mailto:techsweave@gmail.com}{techsweave@gmail.com}\\

    \end{center}
\end{titlepage}


%INDICE
\newpage
\tableofcontents
\newpage

\section{Informazioni riunione}
\begin{itemize}
    \item Data e ora: 2021-05-10 ore 18:00-19:30.
    \item Luogo: Canale Discord del gruppo tramite videochiamata.
    \item Partecipanti interni: Samuele De Simone, Manuel Varo, Tito Scutari, Nicolò Giaccone, Amedeo Meggiolaro, Marco Barbaresco.
    \item Partecipanti esterni: nessuno.
    \item Segretario: Marco Barbaresco.
\end{itemize}
\section{Argomenti affrontati}
Durante l'incontro sono stati discussi il passaggio a un'architettura a microservizi, la pianificazione per il periodo di progettazione di dettaglio ed è stato
richiesto un incontro al proponente.
\section{Verbale della riunione}
\begin{itemize}
    \item Discussione sull'architettura a microservizi;
    \item Brain storming e visualizzazione collettiva di documentazione, guide e tutorial al fine d'individuare i microservizi utili al progetto;
    \item Pianificazione delle attività da svolgere in vista della \textit{Revisione di Qualifica};
    \item Discussione sui requisiti che il prodotto dovrà soddisfare per la \textit{Revisione di Qualifica};
    \item Discussione sul servizio da utilizzare come storage per le immagini;
    \item Discussione su quale framework utilizzare per la realizzazione dell'interfaccia utente (Bootstrap, Material-UI, Tailwind):
          si è deciso di attendere l'incontro con i proponenti e chiedere se hanno delle preferenze;
    \item Suddivisione dei ruoli e dei task per il periodo di progettazione di dettaglio;
    \item Richiesta d'incontro con il proponente \textit{Red Babel}.
\end{itemize}
\newpage
\section{Resoconto}
\begin{center}
    \textbf{\Large Registro delle decisioni}\\
    \vspace{10px}
    \begin{table}[h!]
        \centering
        \rowcolors{2}{logo!10}{logo!40}
        \renewcommand{\arraystretch}{1.8}
        \begin{tabular}{p{200px} p{200px}}
            \rowcolor{logo!70} \textbf{Codice} & \textbf{Decisione}                                                                                                                             \\
            VI\_2021\_05\_10.1                 & Individuati i microservizi utili al progetto, rimandando la decisione sull'effettivo passaggio di architettura alla riunione con i proponenti. \\
            VI\_2021\_05\_10.2                 & Individuati i requisiti che il prodotto dovrà soddisfare in vista della \textit{Revisione di Qualifica}.                                       \\
            VI\_2021\_05\_10.3                 & Assegnati ruoli e task per il periodo di progettazione di dettaglio.                                                                           \\
            VI\_2021\_05\_10.4                 & Utilizzo di Amazon S3 come servizio di storage.                                                                                                \\
            VI\_2021\_05\_10.5                 & Viene richiesto al proponente, tramite piattaforma Slack, un incontro per il 2021-16-05.                                                       \\
        \end{tabular}
    \end{table}
\end{center}

\end{document}
